% ======================================================================
% sections/methods.tex
%
% Methods for the full paper:
%   Triple Humanoid 24/7 Adverse Event Oncology Trial Response Team:
%   4-Site Rotation (v0.7.0)
%
% Nine subsections inside the single \section{Methods} block:
%   1) Code Instruction Layer
%   2) Code Generation Layer (27 src files)
%   3) Schemas and Configs
%   4) Cartesian Coordinate Surface
%   5) Camaraderie Behavior
%   6) Iteration Sweep
%   7) Tests and Error Checks
%   8) Runtime Targets
%   9) Execution
% ======================================================================

\begin{sectionaccent}

\subsection{Code Instruction Layer}
\label{sec:methods-instructions}

The instruction tree at \fp{Clinical-AI-Demos/tree/main/demo-projects/07-humanoid/paper/instructions} carries the project
overview, the commit roadmap, and 10 numbered subdirectories that
describe how the future automated session should author every config,
schema, source file, iteration sweep, language model planner, and
runtime target \cite{kawchak_2026_humanoid_instructions}. Table
\ref{tbl:instruction-tree} maps each subdirectory to its purpose so
the reader can audit the upstream design.

\begin{table}[H]
\captionsetup{type=table}
\caption{Instruction tree subdirectories and their purpose. The 10
numbered directories drive each phase of the code generation. The
\fp{10-repository-update-instructions/} directory is named in the
instructions README repository structure block but is not yet present on disk.}
\label{tbl:instruction-tree}
\small
\begin{tabular}{>{\raggedright\arraybackslash}p{5cm}
                >{\raggedright\arraybackslash}p{10cm}}
\toprule
\textbf{Subdirectory} & \textbf{Purpose} \\
\midrule
\fp{00-project-overview/} & Architecture, swarm, thesis, memory, and camaraderie overview. \\
\fp{01-commit-roadmap/} & Seven commit roadmap with file lists and size budgets. \\
\fp{02-config-instructions/} & Instructions to author the YAML configs in \fp{config/}. \\
\fp{03-schemas-instructions/} & Instructions to author the JSON Schemas in \fp{schemas/}. \\
\fp{04-robotic-instructions/} & Instructions for the dispatcher, the robot loop, and the swarm coordinator. \\
\fp{05-cartesian-instructions/} & Instructions for the kinematics, the sensor to xyz path, and the Cartesian planner. \\
\fp{06-iteration-instructions/} & Instructions for the week runner, the iterate driver, and the Rust runner. \\
\fp{07-llm-planner-instructions/} & Instructions for the language model planner, the broadcaster, the CTCAE grader, and physician escalation. \\
\fp{08-comparison-instructions/} & Instructions for the comparison framework, the report, and the dashboard. \\
\fp{09-runtime-instructions/} & Instructions for the MacOS, Windows, and Linux runtime targets. \\
\fp{10-repository-update-instructions/} & Reserved for the parent README, the releases note, and CHANGELOG updates. \\
\bottomrule
\end{tabular}
\end{table}

\subsection{Code Generation Layer}
\label{sec:methods-codegen}

The generated code tree at \fp{paper/codegen/src/} carries 27 source
files: 24 Python modules, 1 C++ module, and 2 Rust modules. The files
are organized by function rather than by language. Table
\ref{tbl:src-by-function} groups the 27 files into 9 functional
clusters so the reader can navigate by purpose.

\begin{table}[H]
\captionsetup{type=table}
\caption{27 source files in \fp{paper/codegen/src/} grouped by
function. The list is exhaustive; every file in the tree is named.
This grouping highlights that the entire sensor to robot pipeline
fits inside a 27 file tree.}
\label{tbl:src-by-function}
\small
\begin{tabular}{>{\raggedright\arraybackslash}p{4.6cm}
                >{\raggedright\arraybackslash}p{10.4cm}}
\toprule
\textbf{Function} & \textbf{Files} \\
\midrule
Sensor to Cartesian & sensor to xyz, kinematics, cartesian planner, xyz safety. \\
Per site dispatch & h2 dispatcher (3 robot atomic broadcast). \\
Swarm camaraderie & swarm coordinator (191 lines) and peer state tracker. \\
LLM planner and broadcast & llm planner, broadcaster, prompt frozen, comms intellectual, comms physical, cross site bus. \\
CTCAE grading and FDA & ctcae grader, fda rtct submitter (67 lines, ed25519 signed). \\
Physician and audit & physician escalation, check errors. \\
Site runtime and iteration & site runtime, iterate, week runner, ingest, compute. \\
Comparison and figures & compare agent, figures gen. \\
Native code paths & robot loop (C++, 10 Hz), runner and runner main (Rust). \\
\bottomrule
\end{tabular}
\end{table}

The Cartesian pipeline takes raw sensor frames from the H2 EDU
cameras, the infrared beacon, and the ultra wide band peer mesh, and
projects them into a site local x, y, z goal that the Cartesian
planner blends with a quintic profile. The safety guard at
\fp{xyz_safety.py} enforces a per arm cumulative force budget of 15
N during chest compressions, 5 N otherwise, a 1.2 m inter robot
minimum distance at rest, 0.4 m during shared task handoff, and a
22 N cumulative cross robot force during 3 robot patient transfer.
These constants are conservative against the H2 EDU spec sheet, which
is the design intent
\cite{unitree_h2_2026, iec_80601_2_77}.

The per site dispatcher \fp{h2_dispatcher.py} is the single point at
which a language model decision becomes a robot command. The
broadcaster publishes a single atomic message with exactly 3 sub
commands, one for each robot at the site. The dispatcher waits for 3
acknowledgments within 50 milliseconds and extends the SHA-256 audit
chain by one hash. This atomicity is a regulatory feature, not a
performance feature: if any one robot fails to acknowledge, the
entire tick is recorded as failed and the next tick reassigns roles.

The swarm coordinator at 191 lines holds the camaraderie logic. The
compact length is itself evidence that 3 way coordination logic does
not need a peer to peer consensus protocol when a per site
broadcaster pattern is paired with a shared state tracker
\cite{kawchak_2026_humanoid_instructions}. The peer state tracker is a live view of where the other two robots are standing,
what they are doing now and in the future, treating peer
reported sensor readings as first class evidence equal to a robot's sensors.

The language model planner reads the frozen system prompt at
\fp{prompt_frozen.md} and validates each decision against the
\fp{llm_decision.schema.json} schema before broadcast
\cite{anthropic_claude_opus_47_2026}. The broadcaster fans the
validated decision out to all 3 robots at the same 1 Hz tick. The
intellectual communications module carries the language model
exchange over the shared compute fabric. The physical communications
module carries the 60 GHz ultra wide band peer mesh at 200 Hz plus
the infrared beacon listener at 1000 Hz on each robot.

The CTCAE grader applies the grading scale across the 13
adverse event types covered by the decision table
\cite{ctcae_v5_2017}. The FDA RTCT submitter at 67 lines submits
grade 2 or higher records to the RTCT pilot intake API within the 1
hour SLA from HR 9505 \cite{fda_rtct_2026, fda2026realtime}. The 67
line scale is intentional: a small signed JSON submitter has a small
audit surface and the small surface area is a regulatory feature.
Production uses ed25519 signing; simulation uses a placeholder
signature \cite{iec_62304}.

The physician escalation module pages the on call physician on any
grade 3 or higher adverse event and blocks any intravenous
intervention while waiting for a physician acknowledgment. The error
checker traps schema and physics rule violations before broadcast and
provides the audit credibility that the FDA and the IRB require.

The site runtime drives the per site ingestion, schema validation,
broadcast, dispatch, and audit cycle. The week runner orchestrates
the 168 hour 4 site simulation. The iterate driver runs the 32
deterministic iteration sweep across the language model temperature,
the adverse event arrival jitter, the response time variance, and
the camarade aggressiveness axes. The compute module aggregates the
26 key metric vectors. The compare agent generates the natural
language comparison report.

The Rust runner at \fp{runner.rs} and \fp{runner_main.rs} mirrors
the Python iteration driver and runs the 32 sweep against the
deterministic seed contract. The C++ robot loop at
\fp{robot_loop.cpp} drives the 10 Hz in robot motion loop with a 5
millisecond emergency stop reaction. Cross OS compilation is
supported with clang++, g++, and MSVC under their respective C++20
modes.

\subsection{Schemas and Configs}
\label{sec:methods-schemas-configs}

The schemas at \fp{paper/codegen/schemas/} provide 11 JSON Schema
documents at Draft 2020-12 conformance level. Table
\ref{tbl:schemas-configs} lists each schema and config file with its
purpose so the reader can navigate the data layer.

\begin{table}[H]
\captionsetup{type=table}
\caption{Schemas and configs in the codegen tree. Each row points to
a single JSON Schema or YAML config file and names the function it
serves. All 11 schemas validate clean on the sample records held in
\fp{paper/codegen/data/samples.jsonl}.}
\label{tbl:schemas-configs}
\small
\begin{tabular}{>{\raggedright\arraybackslash}p{5cm}
                >{\raggedright\arraybackslash}p{10cm}}
\toprule
\textbf{File} & \textbf{Purpose} \\
\midrule
humanoid command schema & Per robot Cartesian command envelope. \\
swarm message schema & 3 robot atomic broadcast envelope. \\
ae event schema & Sensor sourced adverse event record. \\
ctcae grading schema & CTCAE grading payload. \\
sponsor ackn. schema & Sponsor side acknowledgment. \\
fda rtct submission schema & Signed FDA RTCT intake payload. \\
physician escalation schema & On call pager payload. \\
llm decision schema & 3 sub command output schema. \\
robot camarade state schema & Per robot camaraderie state. \\
peer handoff schema & 2 second hand off envelope. \\
metrics schema & 26 key metric vector. \\
network yaml & PAT-NET-001 network with H2 EDU, Jetson AGX Thor fields. \\
h2 humanoid yaml & 12 H2 EDU robots with dexterous hands and compute module. \\
site coordination yaml & 1 Hz broadcast plus degraded modes. \\
swarm coordination yaml & 4 groups, role priority weights, Hungarian match, 7 invariants. \\
llm loop yaml & Claude Opus 4.7 1M loop at 200 ms tick. \\
escalation rules yaml & CTCAE 3 plus FDA RTCT plus sponsor SLA. \\
iterations yaml & 32 iteration sweep across 5 axes. \\
site frame bounds yaml & Per site Cartesian frame bounds for xyz safety. \\
\bottomrule
\end{tabular}
\end{table}

\subsection{Cartesian Coordinate Surface}
\label{sec:methods-cartesian}

The Cartesian surface is the contract between the language model and
the robot. The language model does not author joint angles. It
authors x, y, z goals expressed in the site local frame. The
kinematics yaml at \fp{paper/codegen/src/kinematics.yaml} encodes the
39 degree of freedom chain for the H2 EDU body, with separate
sections for the dexterous hand (6 fingers, 12 grasp poses, 0.05 N
tactile resolution) and the compute module
\cite{unitree_h2_2026, nvidia_jetson_thor_2026}. The Cartesian
planner at \fp{cartesian_planner.py} blends successive x, y, z goals
with a quintic profile so the robot does not jerk. The xyz safety
guard at \fp{xyz_safety.py} rejects any goal that would violate the
force budget caps, the inter robot minimum distance, or the no fly
zones at the head of the patient bed. Test coverage at
\fp{paper/codegen/tests/test_xyz_safety.py} confirms the bounds
plus the peer envelope plus the no fly zone hold across 6 parametric
test cases (see Section \ref{sec:methods-tests}).

\subsection{Camaraderie Behavior}
\label{sec:methods-camaraderie}

Figure \ref{fig:swarm-dance} carries the verbatim camarade swarm
dance diagram from the codegen tree. The diagram is a top down view
of 3 H2 EDU robots over 10 seconds at 1 Hz frames. Lead stays close
to the patient bed at a 0.4 m offset on the south side. Assist orbits
0.8 m to the west, ready to hand off. Reserve circles a 1.5 m
perimeter, watching the room. At second 6, Lead asks Assist for the
secondary epinephrine; Assist arrives within 2 seconds through the
60 GHz peer mesh. The dexterous H2 EDU hand enables a smooth 0.4 m
hand off without drops.




\begin{figure}[H]
\vspace{-0.15cm}
\centering
\begin{Verbatim}[fontsize=\small]
                       Camarade Swarm Dance (10 second window)
                       =========================================

  Top down view of the patient bed and 3 Unitree H2 EDU robots over 10
  seconds at 1 Hz frames. P is patient bed. L is Lead H2-A. A is Assist
  H2-B. R is Reserve H2-C. Subscripts 0 through 9 are seconds since AE
  peak. Each robot draws its trail.

                                     (north wall)
                       .......................................
                       .                                     .
                       .                R0                   .
                       .                                     .
                       .          R3 R6 R9                   .
                       .                                     .
                       .             A0                      .
                       .         A2  A4  A6   A8             .
                       .                                     .
                       .              L0                     .
                       .              ||                     .
                       .              PP    (patient bed)    .
                       .              ||                     .
                       .              L2                     .
                       .              L4                     .
                       .              L6                     .
                       .              L8                     .
                       .                                     .
                       .......................................
                                    (south wall)

  Observations:
    Lead stays close to patient bed, 0.4 m offset on the south side.
    Assist orbits 0.8 m to the west, ready to hand off.
    Reserve circles a 1.5 m perimeter, watching the room.
    No robot enters another robot 2 second motion envelope.
    At t+6, Lead asks Assist for secondary epinephrine;
    Assist arrives within 2 seconds via 60 GHz UWB hand-off.
    Dexterous H2 EDU hands enable smooth 0.4 m hand-off without drops.
\end{Verbatim}
\caption{Camarade swarm dance at one site over a 10 second adverse event peak window. Lead, Assist, and Reserve trails track the role rotation that the 1 Hz language model broadcast drives.}
\label{fig:swarm-dance}
\end{figure}

\subsection{Iteration Sweep}
\label{sec:methods-iterations}

The iteration sweep at \fp{paper/codegen/src/iterate.py} drives 32
deterministic samples across 5 axes. Figure
\ref{fig:iteration-funnel} carries the verbatim iteration sweep
funnel diagram from the execution tree. Each iteration takes the
iterations yaml config as the source of truth, generates a Latin
Hypercube sample of 32 combinations across the seed, language model
temperature, adverse event arrival jitter, response time variance,
and camarade aggressiveness axes, runs the week runner in fast mode,
and writes one row to \fp{paper/codegen/data/iterations/index.jsonl}.
The DuckDB aggregate at
\fp{paper/execution/data/iterations_aggregate.duckdb} rolls the 32
rows into a 15 column store ready for analytical queries
\cite{duckdb_2026, apache_arrow_2026}.

\begin{figure}[H]
    \hspace{1.6cm}
   \includegraphics[width=0.825\linewidth, trim={0 50 0 0}, clip]{Images/07-06-capability-radar-comparison.png}
    % \vspace{-0.25cm}
    \caption{Four adverse event treatment options compared based on multiple metrics.}
    \label{RadarResponse}
\end{figure}



\begin{figure}[H]
\vspace{-0.15cm}
\centering
\begin{Verbatim}[fontsize=\footnotesize]
                                32 Iteration Sweep + Compare Funnel
                                ===================================
                                config/iterations.yaml
                                +---------------------------+
                                | axes:                     |
                                |   seed: [1..8]            |
                                |   llm_temperature:        |
                                |     [0.0, 0.1, 0.2, 0.3]  |
                                |   ae_arrival_jitter_s:    |
                                |     [0, 30, 60, 120]      |
                                |   response_time_variance: |
                                |     [0.0, 0.05, 0.1, 0.2] |
                                |   camarade_aggressiveness:|
                                |     [0.5, 0.75, 1.0, 1.25]|
                                | iteration_count: 32       |
                                +---------------------------+
                                              |
                                              v
                                +---------------------------+
                                | generate_lhs(n=32)        |
                                | -> 32 Latin Hypercube     |
                                |    sample combinations    |
                                +---------------------------+
                                              |
                                              v
                                +---------------------------+
                                | for each iter in 32:      |
                                |   run_iteration()         |
                                |   -> WeekRunner fast      |
                                |   -> 4 site ProcessPool   |
                                |   -> 3600 ticks per site  |
                                +---------------------------+
                                              |
                                              v
                                +---------------------------+
                                | build_index() ->          |
                                | data/iterations/          |
                                | index.jsonl (32 lines)    |
                                +---------------------------+
                                              |
                                              v
                                +---------------------------+
                                | build_aggregate() ->      |
                                | data/iterations/          |
                                | aggregate.duckdb          |
                                +---------------------------+
                                              |
                                              v
                                +---------------------------+
                                | compute.py 26 key vector  |
                                | vs human_team_baseline    |
                                | weighted score 0.953      |
                                +---------------------------+
                                              |
                                              v
                                +---------------------------+
                                | compare_agent.py          |
                                | natural language report   |
                                +---------------------------+
\end{Verbatim}
\vspace{-0.3cm}
\caption{Iteration sweep funnel showing the source through the Latin Hypercube generator, per iteration week runner, index assembly, DuckDB aggregate, 26 key metric vector, and natural language agent.}
\label{fig:iteration-funnel}
\end{figure}

\subsection{Tests and Error Checks}
\label{sec:methods-tests}

Test coverage at \fp{paper/codegen/tests/} carries 6 pytest modules
that execute 36 tests in 0.22 seconds. The test of xyz safety at
\fp{test_xyz_safety.py} confirms the force budget caps, the inter
robot distance, the cross robot force, and the no fly zone bounds.
The 7 check error scan at \fp{check_errors.py} returns zero issues
across the project. These artifacts provide the audit credibility
that the FDA and other regulators expect for any human factors
review of medical device software life cycle conformance
\cite{iec_62304}.

\subsection{Runtime Targets}
\label{sec:methods-runtime}

The codegen tree runs on MacOS, Windows, and Linux. Table
\ref{tbl:runtime-targets} lists the target platforms and the typical
full iteration wall time. The Dockerfiles at
\fp{paper/codegen/docker/} provide a 4 service per site bundle (the
language model service, the ingest service, the simulator, and the
cross site bus) wired by \fp{docker-compose.yml} at the codegen tree
root.

\begin{table}[H]
\captionsetup{type=table}
\caption{Runtime targets and typical wall time for a full 168 hour
iteration. Reference hardware lines are drawn from
\fp{paper/codegen/README.md} runtime sections.}
\label{tbl:runtime-targets}
\small
\begin{tabular}{>{\raggedright\arraybackslash}p{3.5cm}
                >{\raggedright\arraybackslash}p{6cm}
                >{\raggedright\arraybackslash}p{5.5cm}}
\toprule
\textbf{Operating System} & \textbf{Reference Hardware} & \textbf{Typical Wall Time per Iteration} \\
\midrule
MacOS & Mac Studio M2 Ultra & 12 minutes \\
Windows & Lenovo ThinkStation P8 (96 cores) & 10 minutes \\
Linux & \footnotesize Supermicro AS-2025HS-TNR (192 cores) \normalsize & 8 minutes \\
\bottomrule
\end{tabular}
\end{table}

\subsection{Execution}
\label{sec:methods-execution}

The execution tree at \fp{paper/execution/} carries the 23 per step
logs that capture the runtime outcomes of every executable element
in the codegen tree. The 23 logs walk schema ingest, the 7 check
error scan, the 36 pytest run, the 22 module import set, the per
source file smoke tests for the swarm coordinator, the language
model planner, the CTCAE grader, and the H2 dispatcher, the 4 site
site runtime, the 12 robot dispatcher initialization illustrated in \autoref{RadarResponse}. 

This also includes the week
runner fast mode, the sensor xyz safety smoke, the comms peer smoke,
the FDA physician cross site smoke, the cargo Rust build, the C++
robot loop build under g++ at C++20, the 32 iteration sweep, the 26
metric compute pass, the comparison narrative, the 6 figure
generation pass, the JSON YAML ASCII validation pass, and the final
ruff plus format plus yamllint pass. Figure
\ref{fig:module-dependency} carries the verbatim per site module
dependency diagram at one tick. This diagram is the reader's quickest
path from a single sensor frame to the FDA RTCT submission for that
tick.

\begin{figure}[H]
\vspace{-0.15cm}
\centering
\begin{Verbatim}[fontsize=\footnotesize]
                          Per Site Module Dependency at One Tick (1 Hz)
                         ===============================================
                            
                                       +------------------+
                                       | sensor_to_xyz.py |
                                       | (10 Hz aggregate)|
                                       +--------+---------+
                                                |
                                                v
                                 +--------------+---------------+
                                 | comms_intellectual.py        |
                                 | (pub-sub on Claude Code bus) |
                                 +--------------+---------------+
                                                |
                                                v
                           +--------------------+--------------------+
                           | llm_planner.py 1 Hz, Claude Opus 4.7 1M |
                           | reads prompt_frozen.md system prompt    |
                           | validates llm_decision.schema.json      |
                           +--------------------+--------------------+
                                                |
                                                v (1 broadcast, 3 sub commands)
                           +--------------------+--------------------+
                           | broadcaster.py + h2_dispatcher.py       |
                           | atomic publish, 3 acks/500 ms,          |
                           | extends SHA-256 audit chain             |
                           +-- --+--------------+---------------+----+
                                 |              |               |
                                 v              v               v
                           +--------+--+   +-------+---+   +-------+--+
                           | H2-A Lead |   | H2-B Asst |   | H2-C Rsv |
                           +-----+-----+   +-----+-----+   +---+------+
                                 |               |             |
                                 +---------+-----+-------------+------+
                                           v                   v
                                    +------+------+   +--------+-----+
                                    | comms_      |   | swarm_       |
                                    | physical    |   | coordinator  |
                                    | (UWB + IR)  |   | (roles+inv)  |
                                    +------+------+   +-----+--------+
                                           |                |
                                           +--------+-------+
                                                    v
                                        +-----------+-----------+
                                        | xyz_safety.py guards  |
                                        +-----------+-----------+
                                                    v
                                        +-----------+-----------+
                                        | robot_loop.cpp 10 Hz  |
                                        +-----------+-----------+
                                                    v
                                        +-----------+-----------+
                                        | ctcae_grader.py +     |
                                        | physician_escalation  |
                                        +-----------+-----------+
                                                    v
                                            fda_rtct_submitter
\end{Verbatim}
\caption{Per site module dependency at one 1 Hz tick from sensor frame ingest to FDA RTCT submission. Every module names a single file in \fp{paper/codegen/src/}.}
\label{fig:module-dependency}
\end{figure}

Figure \ref{fig:ae-response-flow} carries the verbatim adverse event
swimlane diagram showing how Lead, Assist, and Reserve coordinate
over a 90 second window from adverse event detection to ready state.

\begin{figure}[H]
\centering
\begin{Verbatim}[fontsize=\footnotesize]
                       AE Response Swimlane (Single AE, 90 seconds)
                       ===============================================

  Time  Lead (H2-A)            Assist (H2-B)            Reserve (H2-C)
  ---   --------------------   ----------------------   ----------------------
  t+0   AE detected, assign    AE detected, prep tools  AE detected, scan area
  t+5   move to bedside        move backup position     stand 1.5 m off side
  t+15  epinephrine ready      hand-off epinephrine     page on-call physician
  t+18  epinephrine injected   pulse oximeter ready     log to audit chain
  t+22  verify spo2 reading    hand-off pulse oximeter  monitor doctor proxi
  t+30  secondary intervention ecg leads attached       fda rtct submission
  t+45  continuous monitoring  vital sign relay         physician arrival
  t+60  transfer to physician  tool inventory restock   sponsor ack log
  t+75  step back, observe     step back, observe       step back, observe
  t+90  ready for next AE      ready for next AE        ready for next AE

  Notes:
    All 3 robots arrive within 90 seconds.
    Hand-offs complete within 2 seconds via 60 GHz UWB.
    Reserve does FDA RTCT submission in parallel with Lead intervention.
    Platform: Unitree H2 EDU. Hands: 6 dexterous fingers, 0.05 N tactile.
    Compute: Jetson AGX Thor 2070 TOPS upgrade path on every robot.
\end{Verbatim}
\caption{Adverse event swimlane over a 90 second window. Lead, Assist, and Reserve trail through detection, tool handoff, intervention, FDA RTCT submission, physician arrival, sponsor acknowledgment log, and ready state in parallel.}
\label{fig:ae-response-flow}
\end{figure}

\subsubsection*{Real life feasibility insight.} Every Cartesian goal
the language model emits is bounded by the force budget caps and the
inter robot minimum distance constants named in the project
instructions \cite{kawchak_2026_humanoid_instructions}. These bounds
are conservative against what the H2 EDU spec sheet permits
\cite{unitree_h2_2026}. The implication is that a future regulator
review can audit the safety envelope from the constants alone,
without rerunning the simulation. The 67 line FDA RTCT submitter and
the 191 line swarm coordinator are intentionally compact for the
same reason: a small audit surface is a regulatory feature
\cite{iec_62304}.


\sectiontable{Methods planning matrix}{%
Instr & M1 & 1 & instructions tree & 10 subdir walk & 1 to 2 pages & Names every subdir & Done \\
Source & M2 & 2 & codegen/src & Per file function & 4 to 5 pages & 191 lines swarm & Done \\
Schemas & M3 & 3 & schemas, config & Schema by schema & 2 pages & Every JSON named & Done \\
Cartes. & M4 & 4 & kinem. plus planner & x, y, z contract & 2 pages & Force caps named & Done \\
Camara. & M5 & 5 & swarm coord. & 191 line walk & 2 pages & Figure 3 ASCII & Done \\
Iter. & M6 & 6 & iterations idx & 32 sweep summary & 1 to 2 pages & 32 across 4 sites & Done \\
Tests & M7 & 7 & check, tests & FDA credibility & 1 to 2 pages & Both named & Done \\
Run. & M8 & 8 & docker plus compose & 3 OS targets & 1 to 2 pages & MacOS Win Linux & Done \\
Exec. & M9 & 9 & exec logs 01 to 23 & Run progression & 2 to 3 pages & Audit trail & Done \\
}

\end{sectionaccent}
